% OpenStax Calculus Volume 3 -- Section 2.5 Exercises: Equations of Lines and Planes in Space
% Exercises reproduced from OpenStax, licensed under CC BY 4.0.
% Source: https://openstax.org/books/calculus-volume-3/pages/2-5-equations-of-lines-and-planes-in-space
\documentclass{exercises}
\title{OpenStax Calculus Volume 3}
\subtitle{Section 2.5 Exercises: Equations of Lines and Planes in Space}
\sourceurl{https://openstax.org/books/calculus-volume-3/pages/2-5-equations-of-lines-and-planes-in-space}
\author{}
\date{}

\begin{document}
\maketitle


In the following exercises, points $P$ and $Q$ are given. Let $L$ be the line passing through points $P$ and $Q$.


(a) Find the vector equation of line $L$.


(b) Find parametric equations of line $L$.


(c) Find symmetric equations of line $L$.


(d) Find parametric equations of the line segment determined by $P$ and $Q$.

\begin{enumerate}[start=1]
\item $P(-3,5,9)$, $Q(4,-7,2)$
\item $P(4,0,5)$, $Q(2,3,1)$
\item $P(-1,0,5)$, $Q(4,0,3)$
\item $P(7,-2,6)$, $Q(-3,0,6)$
\end{enumerate}

For the following exercises, point $P$ and vector $\mathbf{v}$ are given. Let $L$ be the line passing through point $P$ with direction $\mathbf{v}$.


(a) Find parametric equations of line $L$.


(b) Find symmetric equations of line $L$.


(c) Find the intersection of the line with the $xy$-plane.

\begin{enumerate}[resume]
\item $P(1,-2,3)$, $\mathbf{v}=\langle1,2,3\rangle$
\item $P(3,1,5)$, $\mathbf{v}=\langle1,1,1\rangle$
\item $P(3,1,5)$, $\mathbf{v}=\overset{\to}{QR}$, where $Q(2,2,3)$ and $R(3,2,3)$
\item $P(2,3,0)$, $\mathbf{v}=\overset{\to}{QR}$, where $Q(0,4,5)$ and $R(0,4,6)$
\end{enumerate}

For the following exercises, line $L$ is given.


(a) Find point $P$ that belongs to the line and direction vector $\mathbf{v}$ of the line. Express $\mathbf{v}$ in component form.


(b) Find the distance from the origin to line $L$.

\begin{enumerate}[resume]
\item $x=1+t,y=3+t,z=5+4t$, $t\in \mathbb{R}$
\item $-x=y+1,z=2$
\item Find the distance between point $A(-3,1,1)$ and the line of symmetric equations $x=-y=-z$.
\item Find the distance between point $A(4,2,5)$ and the line of parametric equations $x=-1-t,y=-t,z=2$, $t\in \mathbb{R}$.
\end{enumerate}

For the following exercises, lines $L_{1}$ and $L_{2}$ are given.


(a) Verify whether lines $L_{1}$ and $L_{2}$ are parallel.


(b) If the lines $L_{1}$ and $L_{2}$ are parallel, then find the distance between them.

\begin{enumerate}[resume]
\item $L_{1}:x=1+t,y=t,z=2+t$, $t\in \mathbb{R}$, $L_{2}:x-3=y-1=z-3$
\item $L_{1}:x=2,y=1,z=t$, $L_{2}:x=1,y=1,z=2-3t$, $t\in \mathbb{R}$
\item Show that the line passing through points $P(3,1,0)$ and $Q(1,4,-3)$ is perpendicular to the line with equations $x=3+3t,y=1+8t,z=6t$, $t\in \mathbb{R}$.
\item Are the lines of equations $x=-2+2t,y=-6,z=2+6t$ and $x=-1+t,y=1+t,z=t$, $t\in \mathbb{R}$, perpendicular to each other?
\item Find the point of intersection of the lines of equations $x=-2y=3z$ and $x=-5-t,y=-1+t,z=t-11$, $t\in \mathbb{R}$.
\item Find the intersection point of the $x$-axis with the line of parametric equations $x=10+t,y=2-2t,z=-3+3t$, $t\in \mathbb{R}$.
\end{enumerate}

For the following exercises, lines $L_{1}$ and $L_{2}$ are given. Determine whether the lines are equal, parallel but not equal, skew, or intersecting.

\begin{enumerate}[resume]
\item $L_{1}:x=y-1=-z$ and $L_{2}:x-2=-y=\frac{z}{2}$
\item $L_{1}:x=2t,y=0,z=3$, $t\in \mathbb{R}$ and $L_{2}:x=0,y=8+s,z=7+s$, $s\in \mathbb{R}$
\item $L_{1}:x=-1+2t,y=1+3t,z=7t$, $t\in \mathbb{R}$ and $L_{2}:x-1=\frac{2}{3}(y-4)=\frac{2}{7}z-2$
\item $L_{1}:3x=y+1=2z$ and $L_{2}:x=6+2t,y=17+6t,z=9+3t$, $t\in \mathbb{R}$
\item Consider line $L$ of symmetric equations $x-2=-y=\frac{z}{2}$ and point $A(1,1,1)$.
\begin{enumerate}[label=(\alph*)]
  \item Find parametric equations for a line parallel to $L$ that passes through point $A$.
  \item Find symmetric equations of a line skew to $L$ and that passes through point $A$.
  \item Find symmetric equations of a line that intersects $L$ and passes through point $A$.
\end{enumerate}
\item Consider line $L$ of parametric equations $x=t,y=2t,z=3$, $t\in \mathbb{R}$.
\begin{enumerate}[label=(\alph*)]
  \item Find parametric equations for a line parallel to $L$ that passes through the origin.
  \item Find parametric equations of a line skew to $L$ that passes through the origin.
  \item Find symmetric equations of a line that intersects $L$ and passes through the origin.
\end{enumerate}
\end{enumerate}

For the following exercises, point $P$ and vector $\mathbf{n}$ are given.


(a) Find the scalar equation of the plane that passes through $P$ and has normal vector $\mathbf{n}$.


(b) Find the general form of the equation of the plane that passes through $P$ and has normal vector $\mathbf{n}$.

\begin{enumerate}[resume]
\item $P(0,0,0)$, $\mathbf{n}=3\mathbf{i}-2\mathbf{j}+4\mathbf{k}$
\item $P(3,2,2)$, $\mathbf{n}=2\mathbf{i}+3\mathbf{j}-\mathbf{k}$
\item $P(1,2,3)$, $\mathbf{n}=\langle1,2,3\rangle$
\item $P(0,0,0)$, $\mathbf{n}=\langle-3,2,-1\rangle$
\end{enumerate}

For the following exercises, the equation of a plane is given.


(a) Find normal vector $\mathbf{n}$ to the plane. Express $\mathbf{n}$ using standard unit vectors.


(b) Find the intersections of the plane with the coordinate axes.


(c) Sketch the plane.

\begin{enumerate}[resume]
\item \techrequired $4x+5y+10z-20=0$
\item $3x+4y-12=0$
\item $3x-2y+4z=0$
\item $x+z=0$
\item Given point $P(1,2,3)$ and vector $\mathbf{n}=\mathbf{i}+\mathbf{j}$, find point $Q$ on the $x$-axis such that $\overset{\to}{PQ}$ and $\mathbf{n}$ are orthogonal.
\item Show there is no plane perpendicular to $\mathbf{n}=\mathbf{i}+\mathbf{j}$ that passes through points $P(1,2,3)$ and $Q(2,3,4)$.
\item Find parametric equations of the line passing through point $P(-2,1,3)$ that is perpendicular to the plane of equation $2x-3y+z=7$.
\item Find symmetric equations of the line passing through point $P(2,5,4)$ that is perpendicular to the plane of equation $2x+3y-5z=0$.
\item Show that line $\frac{x-1}{2}=\frac{y+1}{3}=\frac{z-2}{4}$ is parallel to plane $x-2y+z=6$.
\item Find the real number $\alpha$ such that the line of parametric equations $x=t,y=2-t,z=3+t$, $t\in \mathbb{R}$ is parallel to the plane of equation $\alpha x+5y+z-10=0$.
\end{enumerate}

For the following exercises, points $P$, $Q$, and $R$ are given.


(a) Find the general equation of the plane passing through $P$, $Q$, and $R$.


(b) Write the vector equation $\mathbf{n}\cdot \overset{\to}{PS}=0$ of the plane at a., where $S(x,y,z)$ is an arbitrary point of the plane.


(c) Find parametric equations of the line passing through the origin that is perpendicular to the plane passing through $P$, $Q$, and $R$.

\begin{enumerate}[resume]
\item $P(1,1,1),Q(2,4,3)$, and $R(-1,-2,-1)$
\item $P(-2,1,4),Q(3,1,3)$, and $R(-2,1,0)$
\item Consider the planes of equations $x+y+z=1$ and $x+z=0$.
\begin{enumerate}[label=(\alph*)]
  \item Show that the planes intersect.
  \item Find symmetric equations of the line passing through point $P(1,4,6)$ that is parallel to the line of intersection of the planes.
\end{enumerate}
\item Consider the planes of equations $-y+z-2=0$ and $x-y=0$.
\begin{enumerate}[label=(\alph*)]
  \item Show that the planes intersect.
  \item Find parametric equations of the line passing through point $P(-8,0,2)$ that is parallel to the line of intersection of the planes.
\end{enumerate}
\item Find the scalar equation of the plane that passes through point $P(-1,2,1)$ and is perpendicular to the line of intersection of planes $x+y-z-2=0$ and $2x-y+3z-1=0$.
\item Find the general equation of the plane that passes through the origin and is perpendicular to the line of intersection of planes $-x+y+2=0$ and $z-3=0$.
\item Determine whether the line of parametric equations $x=1+2t,y=-2t,z=2+t$, $t\in \mathbb{R}$ intersects the plane with equation $3x+4y+6z-7=0$. If it does intersect, find the point of intersection.
\item Determine whether the line of parametric equations $x=5,y=4-t,z=2t$, $t\in \mathbb{R}$ intersects the plane with equation $2x-y+z=5$. If it does intersect, find the point of intersection.
\item Find the distance from point $P(1,5,-4)$ to the plane of equation $3x-y+2z-6=0$.
\item Find the distance from point $P(1,-2,3)$ to the plane of equation $(x-3)+2(y+1)-4z=0$.
\end{enumerate}

For the following exercises, the equations of two planes are given.


(a) Determine whether the planes are parallel, orthogonal, or neither.


(b) If the planes are neither parallel nor orthogonal, then find the measure of the angle between the planes. Express the answer in degrees rounded to the nearest integer.

\begin{enumerate}[resume]
\item \techrequired $x+y+z=0$, $2x-y+z-7=0$
\item $5x-3y+z=4$, $x+4y+7z=1$
\item $x-5y-z=1$, $5x-25y-5z=-3$
\item \techrequired $x-3y+6z=4$, $5x+y-z=4$
\item Show that the lines of equations $x=t,y=1+t,z=2+t$, $t\in \mathbb{R}$, and $\frac{x}{2}=\frac{y-1}{3}=z-3$ are skew, and find the distance between them.
\item Show that the lines of equations $x=-1+t,y=-2+t,z=3t$, $t\in \mathbb{R}$, and $x=5+s,y=-8+2s,z=7s$, $s\in \mathbb{R}$ are skew, and find the distance between them.
\item Consider point $C(-3,2,4)$ and the plane of equation $2x+4y-3z=8$.
\begin{enumerate}[label=(\alph*)]
  \item Find the radius of the sphere with center $C$ tangent to the given plane.
  \item Find point $P$ of tangency.
\end{enumerate}
\item Consider the plane of equation $x-y-z-8=0$.
\begin{enumerate}[label=(\alph*)]
  \item Find an equation of the sphere with center $C$ at the origin that is tangent to the given plane.
  \item Find parametric equations of the line passing through the origin and the point of tangency.
\end{enumerate}
\item Two children are playing with a ball. The girl throws the ball to the boy. The ball travels in the air, curves $3$ ft to the right, and falls $5$ ft away from the girl (see the following figure). If the plane that contains the trajectory of the ball is perpendicular to the ground, find its equation. \\ \textit{[Figure: This figure is the image of two children throwing a ball. The path of the ball is represented with an arc. The distance from the child throwing the ball to the point where the ball hits is 5 feet. The distance from the second child to where the ball hits is 3 feet.]}
\item \techrequired John allocates $d$ dollars to consume monthly three goods of prices $a$, $b$, and $c$. In this context, the budget equation is defined as $ax+by+cz=d$, where $x\ge0,y\ge0$, and $z\ge0$ represent the number of items bought from each of the goods. The budget set is given by $\{(x,y,z)|ax+by+cz\le d,x\ge0,y\ge0,z\ge0\}$, and the budget plane is the part of the plane of equation $ax+by+cz=d$ for which $x\ge0,y\ge0$, and $z\ge0$. Consider $a=\$8$, $b=\$5$, $c=\$10$, and $d=\$500$.
\begin{enumerate}[label=(\alph*)]
  \item Use a CAS to graph the budget set and budget plane.
  \item For $z=25$, find the new budget equation and graph the budget set in the same system of coordinates.
\end{enumerate}
\item \techrequired Consider $\mathbf{r}(t)=\langle \sin t,\cos t,2t\rangle$ the position vector of a particle at time $t\in[0,3]$, where the components of $\mathbf{r}$ are expressed in centimeters and time is measured in seconds. Let $\overset{\to}{OP}$ be the position vector of the particle after $1$ sec.
\begin{enumerate}[label=(\alph*)]
  \item Determine the velocity vector $\mathbf{v}(1)$ of the particle after $1$ sec.
  \item Find the scalar equation of the plane that is perpendicular to $\mathbf{v}(1)$ and passes through point $P$. This plane is called the normal plane to the path of the particle at point $P$.
  \item Use a CAS to visualize the path of the particle along with the velocity vector and normal plane at point $P$.
\end{enumerate}
\item \techrequired A solar panel is mounted on the roof of a house. The panel may be regarded as positioned at the points of coordinates (in meters) $A(8,0,0)$, $B(8,18,0)$, $C(0,18,8)$, and $D(0,0,8)$ (see the following figure). \\ \textit{[Figure: This figure is a picture of a rectangular solar grid on a roof. The corners of the rectangle are labeled A, B, C, D. There are two vectors, the first is from A to D. The second is from A to B.]}
\begin{enumerate}[label=(\alph*)]
  \item Find the general form of the equation of the plane that contains the solar panel by using points $A$, $B$, and $C$, and show that its normal vector is equivalent to $\overset{\to}{AB}\times \overset{\to}{AD}$.
  \item Find parametric equations of line $L_{1}$ that passes through the center of the solar panel and has direction vector $\mathbf{s}=\frac{1}{\sqrt{3}}\mathbf{i}+\frac{1}{\sqrt{3}}\mathbf{j}+\frac{1}{\sqrt{3}}\mathbf{k}$, which points toward the position of the Sun at a particular time of day.
  \item Find symmetric equations of line $L_{2}$ that passes through the center of the solar panel and is perpendicular to it.
  \item Determine the angle of elevation of the Sun above the solar panel by using the angle between lines $L_{1}$ and $L_{2}$.
\end{enumerate}
\end{enumerate}

\end{document}
